% ============================================================
% TUMOR BIOLOGY CONTRIBUTIONS TO PAPER
% Author: Tumor Biology Agent (agent-mnez9frj)
% Sections: Results (biological validation), Discussion (clinical translation + pDC)
% ============================================================

% --- RESULTS: Biological validation ---
% Insert into Section 2.X "Biological validation"

\subsection{Biological validation: stepwise elimination of an immunosuppressive Treg community}

To validate the biological significance of subpopulation disruption, we traced the fate of T cell Cluster~5, a skin-cancer-specific GITR$^+$ activated regulatory T cell (Treg) functional community (1,464 cells; GITR=69.8\%, TIGIT=65\%, FOXP3=28\%; 86\% from SSCC/skin normal samples). This cluster contains 266 cells co-expressing the triple combination FOXP3$^+$GITR$^+$TIGIT$^+$CCR8$^-$, representing a previously uncharacterized Treg subtype.

\textbf{Full-pipeline fate tracking.} We quantified the cumulative impact of each pipeline step on Cluster~5 (Fig.~5a):
\begin{itemize}
    \item \textbf{QC/filtering}: Activated Tregs exhibit elevated mitochondrial gene fractions due to high metabolic activity (oxidative phosphorylation). Standard mt\% thresholds ($>$20\%) disproportionately remove these cells---a bias not previously recognized in the QC literature.
    \item \textbf{HVG selection}: The key marker TNFRSF18 (GITR), expressed in $<$1\% of total cells, was excluded from the global HVG set despite being the top distinguishing gene for this subpopulation.
    \item \textbf{PCA compression}: With 50 principal components optimized for global variance, the directions separating GITR$^+$ Tregs from mainstream Tregs received negligible weight, compressing the subpopulation to a narrow region indistinguishable from neighboring CD4$^+$ T cell populations.
    \item \textbf{Batch integration}: With 86\% of cells originating from a single cancer type (SSCC), cross-batch alignment forces these cells toward mainstream Tregs in other batches. Survival rate after Harmony integration: 24\%, with cells dispersed across 8 post-integration clusters (Fig.~5b).
\end{itemize}

The cumulative effect is complete elimination: a biologically distinct immunosuppressive community is reduced from a coherent functional unit to undetectable fragments scattered across unrelated clusters. RASI rescues this subpopulation (survival: 24\% $\to$ 84\%) through batch-diversity weighting that recognizes the tissue-specific nature of this community (Fig.~5c).

\textbf{Independent wet-lab validation.} Independently conducted experiments (Methods) confirmed the functional significance of TIGIT$^+$CCR8$^-$ Tregs isolated from this computationally identified subpopulation. In co-culture assays with CD8$^+$ T cells (Treg:T = 1:4, anti-CD3/CD28 stimulation, 72h), TIGIT$^+$CCR8$^-$ Tregs exhibited potent immunosuppression across three complementary readouts (Fig.~5d--f):
\begin{itemize}
    \item T cell activation (CD69$^+$): 20.5\% $\to$ 5.0\% ($\sim$4.1$\times$ reduction, $p<0.001$)
    \item Cytokine secretion (IFN-$\gamma$ ELISA): 305 $\to$ 115~pg/mL ($\sim$2.7$\times$ reduction, $p<0.001$)
    \item Tumor killing (LDH release, A549 target cells): 63\% $\to$ 19\% ($\sim$3.3$\times$ reduction, $p<0.001$)
\end{itemize}

The novelty of this subpopulation lies in the TIGIT$^+$CCR8$^-$ combination: while CCR8$^+$ tumor-infiltrating Tregs are well characterized as a tumor-specific immunosuppressive population\cite{declerck2018}, TIGIT$^+$CCR8$^-$ Tregs represent an alternative immunosuppressive pathway operating through the TIGIT--CD155 axis rather than the canonical CCR8 axis. Their destruction during standard integration means this alternative immunosuppressive mechanism would be invisible in any downstream analysis of integrated data.


% --- DISCUSSION: Clinical translation ---
% Insert into Discussion section

\subsection*{Clinical implications of rare subpopulation loss}

The elimination of rare immunosuppressive communities during batch integration has direct consequences for precision oncology. The GITR$^+$TIGIT$^+$CCR8$^-$ Treg subpopulation identified in our study exemplifies a class of therapeutically actionable rare populations that current workflows systematically destroy.

GITR (TNFRSF18) and TIGIT are both active immunotherapy targets. Anti-GITR antibodies (TRX518, MK-4166, BMS-986156) are in Phase~I/II clinical trials, acting through dual mechanisms: depleting GITR$^+$ Tregs via ADCC and co-stimulating effector T cells. Anti-TIGIT antibodies (tiragolumab) have advanced to Phase~III trials across multiple tumor types. The co-expression of both targets on a single Treg subpopulation, combined with the absence of CCR8, suggests a rationale for anti-GITR + anti-TIGIT combination therapy specifically in skin cancers---a hypothesis that could not be generated from standard integrated atlases where this subpopulation has been eliminated.

More broadly, any cell--cell communication analysis (CellChat, CellPhoneDB) performed on integrated data will miss tissue-specific immune interactions mediated by subpopulations destroyed during integration. Our finding that 42\% of T cell subpopulations and 38\% of mast cell subpopulations are disrupted by standard Harmony integration implies that a substantial fraction of the tumor immune microenvironment is misrepresented in current integrated atlases.

\subsection*{Scale-dependent elimination: implications for atlas-scale studies}

The most striking biological finding is the scale dependence of rare subpopulation loss. pDC---a well-characterized cell type with distinct marker genes (IRF7, TLR7/9, IL3RA, CLEC4C)---was fully preserved when integrating 8 batches ($\sim$105k cells; 99\% recall) but severely disrupted at 103 batches ($\sim$2.26M cells; NP = 0.231, lowest across all cell types).

This reveals a critical implication for the field's trajectory toward ever-larger atlas projects: \textbf{methods benchmarked on small datasets may appear to preserve rare subpopulations while failing catastrophically at atlas scale}. The mechanism is a cascade amplification: as batch count increases, (i) batch-specific HVGs crowd out rare-type markers, (ii) batch variance dominates PCA directions over rare-type variance, (iii) kNN graphs accumulate spurious cross-type connections, and (iv) integration algorithms find false MNN anchors between rare cells and abundant types. Each step amplifies the preceding bias, producing a nonlinear ``phase transition'' in rare subpopulation survival as dataset scale increases (Fig.~4).

pDC play dual roles in tumor immunity---both anti-tumor (IFN-I production) and pro-tumor (immune tolerance via IDO/ICOS-L). Different tumor types harbor functionally distinct pDC states, and the internal heterogeneity of pDC (at least 26 subclusters in our pre-integration analysis, including NHL-specific Cluster~24 with only 93 cells) is entirely obliterated after large-scale integration. This means that pan-cancer studies using standard integrated data will incorrectly conclude that pDC are transcriptionally homogeneous across cancer types---a finding that is an artifact of integration, not biology.
